% ======================================================================
%  LATEX LAB-EXAM COMPLETE REFERENCE
%  ------------------------------------------------------------------
%  ONE self-contained, printable manual. Compile with pdflatex TWICE
%  (so the Table of Contents and cross-references resolve).
%
%    pdflatex latex_exam_reference.tex
%    pdflatex latex_exam_reference.tex
%
%  Every symbol is shown beside the code that makes it. Every table,
%  figure, list and math construct is shown as CODE + RENDERED OUTPUT.
%  Read it, compile it, print it, take it to the lab.
% ======================================================================
\documentclass[11pt]{article}

% ---------- Page ----------
\usepackage[a4paper,margin=1in]{geometry}
\usepackage[T1]{fontenc}
\usepackage[utf8]{inputenc}
\usepackage{textcomp}

% ---------- Colour (BEFORE tcolorbox / tikz) ----------
\usepackage[table,dvipsnames]{xcolor}

% ---------- Mathematics ----------
\usepackage{mathtools}   % superset of amsmath
\usepackage{amssymb}
\usepackage{amsfonts}
\usepackage{amsthm}
\usepackage{mathrsfs}    % \mathscr
\usepackage{bm}          % \bm bold math

% ---------- Graphics ----------
\usepackage{graphicx}
\usepackage{mwe}         % supplies example-image, example-image-a/b/c

% ---------- Tables ----------
\usepackage{array}
\usepackage{multirow}
\usepackage{makecell}
\usepackage{booktabs}
\usepackage{tabularx}

% ---------- Floats ----------
\usepackage{subcaption}
\usepackage{wrapfig}
\usepackage{rotating}
\usepackage{float}

% ---------- Lists / code / authors / boxes ----------
\usepackage{enumitem}
\usepackage{listings}
\usepackage{authblk}
\usepackage[most]{tcolorbox}
\usepackage{pifont}      % \ding{51} check, \ding{55} cross

% ---------- Hyperlinks (near last) ----------
\usepackage{hyperref}
\hypersetup{colorlinks=true, linkcolor=RoyalBlue, urlcolor=magenta,
            citecolor=teal, filecolor=RoyalBlue}

% ---------- Algorithms (after hyperref, per algorithm2e docs) ----------
\usepackage[ruled,vlined,linesnumbered]{algorithm2e}

% ======================================================================
%  HELPER MACROS
% ======================================================================
% \sym{<cmd>} : renders the symbol AND prints its literal code side by side
% Used inside tabulars: \sym{\alpha} & \sym{\beta} & \sym{\gamma} \\
\newcommand{\sym}[1]{\ensuremath{#1} & {\ttfamily\footnotesize\detokenize{#1}}}
% text-mode version (for \dag, \S, \%, ...)
\newcommand{\tsym}[1]{#1 & {\ttfamily\footnotesize\detokenize{#1}}}

% a light box to show "code -> output" demos
\newtcolorbox{demo}{colback=black!3, colframe=black!25, boxrule=0.4pt,
  left=4pt, right=4pt, top=3pt, bottom=3pt, breakable}

% listing styles
\lstdefinestyle{tex}{
  language=[LaTeX]TeX, basicstyle=\ttfamily\small,
  keywordstyle=\color{RoyalBlue}\bfseries,
  commentstyle=\color{gray}\itshape,
  backgroundcolor=\color{black!3}, frame=single, framerule=0.3pt,
  breaklines=true, columns=fullflexible, showstringspaces=false,
  keepspaces=true, morekeywords={multirow,multicolumn,rowcolor,cellcolor,
  includegraphics,subfigure,caption,label,cline,toprule,midrule,bottomrule,
  href,url,cite,eqref,ref,dfrac,align,cases,bmatrix,pmatrix}}
\lstdefinestyle{py}{
  language=Python, basicstyle=\ttfamily\small,
  keywordstyle=\color{blue}\bfseries,
  commentstyle=\color{OliveGreen}\itshape, stringstyle=\color{BrickRed},
  numbers=left, numberstyle=\tiny\color{gray}, numbersep=8pt,
  frame=single, breaklines=true, showstringspaces=false}

% theorem environments
\theoremstyle{plain}
\newtheorem{theorem}{Theorem}[section]
\newtheorem{proposition}[theorem]{Proposition}
\newtheorem{lemma}[theorem]{Lemma}
\theoremstyle{definition}
\newtheorem{definition}{Definition}[section]
\theoremstyle{remark}
\newtheorem*{remark}{Remark}

% coloured theorem box
\newtcbtheorem[number within=section]{keyresult}{Key Result}
  {colback=yellow!8, colframe=orange!75!black, fonttitle=\bfseries}{kr}

\setlength{\parskip}{4pt}
\setlength{\parindent}{0pt}

% ======================================================================
\title{\textbf{\LaTeX{} Lab-Exam Complete Reference}\\[2pt]
       \large Syntax, Symbols, Tables, Figures \& Explanations --- One Sheet}
\author{Personal exam documentation}
\date{\today}
% ======================================================================

\begin{document}
\maketitle
\begin{center}
\fbox{\parbox{0.9\textwidth}{\centering\small
Compile \textbf{twice} with \texttt{pdflatex}. Every symbol appears beside
its code. Every construct is shown as \textbf{code} then \textbf{output}.
Ctrl--F any feature during the exam.}}
\end{center}
\tableofcontents
\newpage

% ======================================================================
\section{The 30-Second Preamble}
Everything starts with a \emph{document class} and a \emph{preamble} (the
zone before \verb|\begin{document}| where you load packages). Copy one of
these first, then start typing.

\textbf{Minimal} --- math, images, colour, links:
\begin{lstlisting}[style=tex]
\documentclass[11pt]{article}
\usepackage[margin=1in]{geometry}
\usepackage{amsmath,amssymb,graphicx,xcolor,hyperref}
\begin{document}
% ... your content ...
\end{document}
\end{lstlisting}

\textbf{Exam-standard} --- add only what the paper needs:
\begin{lstlisting}[style=tex]
\documentclass[11pt]{article}
\usepackage[a4paper,margin=1in]{geometry}
\usepackage[T1]{fontenc}\usepackage[utf8]{inputenc}
\usepackage{amsmath,amssymb,amsthm}   % maths + theorems
\usepackage{graphicx}                 % images
\usepackage[table,dvipsnames]{xcolor} % colour + coloured tables
\usepackage{multirow,array,booktabs}  % strong tables
\usepackage{subcaption,wrapfig}       % subfigures, wrapped figures
\usepackage{enumitem}                 % custom list labels
\usepackage{listings}                 % code
\usepackage{authblk}                  % numbered author affiliations
\usepackage{hyperref}                 % LOAD LAST: links + clickable TOC
\end{lstlisting}

\begin{tcolorbox}[colback=RoyalBlue!6,colframe=RoyalBlue!60,title=Golden rules]
\ding{51}\ \texttt{amsmath} for equations, \texttt{graphicx} for images,
\texttt{xcolor} for colour, \texttt{hyperref} \textbf{last}.\\
\ding{51}\ A \emph{package} is loaded with \verb|\usepackage{name}|; options
go in brackets: \verb|\usepackage[opt]{name}|.\\
\ding{51}\ Comment out a line with \texttt{\%}. Escape specials:
\verb|\% \$ \& \# \_ \{ \}|.
\end{tcolorbox}

% ======================================================================
\section{Document Structure}
\subsection{Title block and authors}
The \texttt{authblk} package wires each author to a numbered affiliation:
\begin{lstlisting}[style=tex]
\title{My Paper Title}
\author[1]{Alice Smith}\author[2]{Bob Jones}
\affil[1]{University A, City}
\affil[2]{Institute B, City}
\date{}                 % empty {} deletes the date; omit it -> today's date
\maketitle              % first line INSIDE \begin{document}
\end{lstlisting}

\subsection{Sectioning}
The commands, from largest to smallest, are
\verb|\section|, \verb|\subsection|, \verb|\subsubsection|,
\verb|\paragraph|, \verb|\subparagraph|. Each auto-numbers and enters the
Table of Contents. Add a star (\verb|\section*{...}|) to make a heading
\emph{unnumbered} and keep it \emph{out} of the TOC --- useful for
``Abstract'' or an unnumbered ``Circuit Model'' heading.

\subsection{Table of contents}
\verb|\tableofcontents| prints every heading with page numbers, and becomes
clickable once \texttt{hyperref} is loaded. Because it reads data recorded
on the \emph{previous} run, you must compile \textbf{twice}.

% ======================================================================
\section{Text Formatting, Fonts, Sizes and Colour}
\subsection{Weights and shapes}
\begin{tabular}{@{}ll@{\hspace{2em}}ll@{}}
\toprule
Output & Code & Output & Code\\\midrule
\textbf{bold} & \verb|\textbf{..}| & \textit{italic} & \verb|\textit{..}|\\
\underline{under} & \verb|\underline{..}| & \texttt{mono} & \verb|\texttt{..}|\\
\textsc{Small Caps} & \verb|\textsc{..}| & \emph{emphasis} & \verb|\emph{..}|\\
\textsf{sans serif} & \verb|\textsf{..}| & \textrm{roman} & \verb|\textrm{..}|\\
\bottomrule
\end{tabular}

\subsection{Size ladder}
From tiny to huge (declared inside braces to keep the scope local):
{\tiny tiny} {\scriptsize script} {\footnotesize footnote} {\small small}
{\normalsize normal} {\large large} {\Large Large} {\LARGE LARGE}
{\huge huge}. Codes: \verb|\tiny \scriptsize \footnotesize \small|
\verb|\normalsize \large \Large \LARGE \huge \Huge|.

\subsection{Colour}
\begin{demo}
\ttfamily\small
\{\textbackslash color\{red\}text\} $\to$ {\color{red}text}\quad
\textbackslash textcolor\{blue\}\{text\} $\to$ \textcolor{blue}{text}\\
\textbackslash colorbox\{yellow\}\{hi\} $\to$ \colorbox{yellow}{hi}\quad
\textbackslash fcolorbox\{black\}\{green!20\}\{hi\} $\to$
  \fcolorbox{black}{green!20}{hi}
\end{demo}
Define your own colour once, reuse it anywhere:
\verb|\definecolor{myblue}{RGB}{0,80,180}| then
\verb|\textcolor{myblue}{...}|. The \verb|!| operator mixes:
\verb|blue!30| means 30\,\% blue, 70\,\% white.

% ======================================================================
\section{Lists and Nested Lists}
Three base environments: \texttt{itemize} (bullets), \texttt{enumerate}
(numbers), \texttt{description} (bold terms). The \texttt{enumitem} package
lets you set the label with \verb|[label=...]| and reduce spacing with
\verb|[noitemsep]|.

\begin{demo}
\begin{minipage}[t]{0.55\textwidth}
\begin{lstlisting}[style=tex]
\begin{enumerate}[label=(\alph*)]
  \item first
  \item second
    \begin{itemize}[label=$\star$]
      \item star bullet
      \item[$\oplus$] one-off marker
    \end{itemize}
\end{enumerate}
\end{lstlisting}
\end{minipage}\hfill
\begin{minipage}[t]{0.4\textwidth}
\begin{enumerate}[label=(\alph*)]
  \item first
  \item second
    \begin{itemize}[label=$\star$]
      \item star bullet
      \item[$\oplus$] one-off marker
    \end{itemize}
\end{enumerate}
\end{minipage}
\end{demo}

Common label patterns:
\begin{tabular}{@{}ll@{\hspace{2em}}ll@{}}
\toprule
Produces & \texttt{label=} & Produces & \texttt{label=}\\\midrule
1. 2. 3.       & \verb|\arabic*.|   & a) b) c)   & \verb|\alph*)|\\
(a) (b)        & \verb|(\alph*)|    & I. II. III.& \verb|\Roman*.|\\
i. ii. iii.    & \verb|\roman*.|    & A) B)      & \verb|\Alph*)|\\
$\dagger$      & \verb|$\dagger$|   & $\triangleright$ & \verb|$\triangleright$|\\
\bottomrule
\end{tabular}

The \texttt{description} list (definition style):
\begin{demo}
\begin{minipage}[t]{0.5\textwidth}
\begin{lstlisting}[style=tex]
\begin{description}
  \item[Term A] its meaning.
  \item[Term B] its meaning.
\end{description}
\end{lstlisting}
\end{minipage}\hfill
\begin{minipage}[t]{0.45\textwidth}
\begin{description}
  \item[Term A] its meaning.
  \item[Term B] its meaning.
\end{description}
\end{minipage}
\end{demo}
Useful keys: \verb|resume| (continue a previous list's count),
\verb|start=3| (begin at 3), \verb|itemsep=2pt|, \verb|leftmargin=1cm|.

% ======================================================================
\section{Mathematics}
\subsection{The two modes}
\emph{Inline} math sits in the running line with \verb|$...$|: the mass
energy relation $E = mc^2$. \emph{Display} math is centred on its own line
with \verb|\[...\]| (never the obsolete \verb|$$...$$|):
\[ \int_0^\infty e^{-x}\,dx = 1. \]

\subsection{Numbered and multi-line equations}
\begin{tabular}{@{}p{0.32\textwidth}p{0.63\textwidth}@{}}
\toprule
Environment & Purpose\\\midrule
\texttt{equation} & one numbered line\\
\texttt{align} & many lines, \verb|&| = alignment column, each numbered\\
\texttt{align*} & same but no numbers\\
\texttt{aligned} & align inside ONE \texttt{equation} (single number)\\
\texttt{gather} & each line centred with its own number (no \verb|&|)\\
\texttt{multline} & one long formula broken across lines\\
\texttt{cases} & piecewise definitions\\
\texttt{subequations} & wraps lines to number them (1a),(1b)\\
\bottomrule
\end{tabular}

\begin{demo}
\begin{minipage}[t]{0.55\textwidth}
\begin{lstlisting}[style=tex]
\begin{align}
  f(x) &= (x+1)^2      \\
       &= x^2 + 2x + 1
\end{align}
\end{lstlisting}
\end{minipage}\hfill
\begin{minipage}[t]{0.4\textwidth}
\begin{align*}
  f(x) &= (x+1)^2 \\
       &= x^2 + 2x + 1
\end{align*}
\end{minipage}
\end{demo}
Suppress one line's number with \verb|\nonumber|; force a custom number with
\verb|\tag{$\ast$}|. The \verb|&| always marks where lines align; \verb|\\|
ends a line.

\subsection{Fractions, roots and big operators}
\begin{tabular}{@{}lll@{\hspace{1.5em}}lll@{}}
\toprule
Out & Code & & Out & Code &\\\midrule
$\dfrac{a}{b}$ & \verb|\dfrac{a}{b}| & & $\sqrt{x}$ & \verb|\sqrt{x}| &\\
$\tfrac{a}{b}$ & \verb|\tfrac{a}{b}| & & $\sqrt[n]{x}$ & \verb|\sqrt[n]{x}| &\\
$\binom{n}{k}$ & \verb|\binom{n}{k}| & & $x^{2n}$ & \verb|x^{2n}| &\\
$\displaystyle\sum_{i=1}^{n}$ & \verb|\sum_{i=1}^{n}| & &
$\displaystyle\int_a^b$ & \verb|\int_a^b| &\\
$\displaystyle\prod_{k}$ & \verb|\prod_{k}| & &
$\displaystyle\lim_{x\to\infty}$ & \verb|\lim_{x\to\infty}| &\\
\bottomrule
\end{tabular}

\subsection{Subscripts, superscripts and deep nesting}
Always \textbf{group with braces}. \verb|e^-st| breaks; \verb|e^{-st}|
works. Nest freely: \verb|x_{i_{j_{k}}}| gives $x_{i_{j_{k}}}$;
\verb|e^{-t/(RC)}| gives $e^{-t/(RC)}$; a symbol with both indices,
\verb|\Gamma^{\mu}_{\alpha\beta}| gives $\Gamma^{\mu}_{\alpha\beta}$; a
tensor with a gap, \verb|A^{\mu}_{\ \nu}| gives $A^{\mu}_{\ \nu}$. A
stacked multi-line subscript uses \verb|\substack|:
\[ \sum_{\substack{i=1\\ i\ne j}}^{n} a_{ij}. \]

\subsection{Delimiters that grow}
\verb|\left...\right| auto-sizes to the content; they must be paired
(use \verb|\right.| for an invisible partner).
\begin{demo}
\ttfamily\small
\textbackslash left(\textbackslash dfrac\{a\}\{b\}\textbackslash right)
$\to \left(\dfrac{a}{b}\right)$\quad
\textbackslash left[\ \textbackslash right]$\to\left[\ \rule{0pt}{2ex}\right]$\quad
\textbackslash left\textbackslash\{\ \textbackslash right\textbackslash\}
$\to\left\{\ \rule{0pt}{2ex}\right\}$\quad
\textbackslash left|\ \textbackslash right|$\to\left|\ \rule{0pt}{2ex}\right|$\quad
\textbackslash langle\ \textbackslash rangle $\to\langle\ \rangle$
\end{demo}
Manual sizes (small$\to$big): \verb|\big \Big \bigg \Bigg|. Floor/ceiling:
\verb|\lfloor x \rfloor| $=\lfloor x\rfloor$,
\verb|\lceil x \rceil| $=\lceil x\rceil$.

\subsection{Matrices and cases}
\begin{demo}
\begin{minipage}[t]{0.55\textwidth}
\begin{lstlisting}[style=tex]
\begin{pmatrix} a & b\\ c & d \end{pmatrix}
\begin{bmatrix} 1 & 0\\ 0 & 1 \end{bmatrix}
f(x)=\begin{cases}
   x^2 & x\ge 0\\ -x & x<0 \end{cases}
\end{lstlisting}
\end{minipage}\hfill
\begin{minipage}[t]{0.4\textwidth}
\[ \begin{pmatrix} a & b\\ c & d \end{pmatrix}
   \begin{bmatrix} 1 & 0\\ 0 & 1 \end{bmatrix} \]
\[ f(x)=\begin{cases} x^2 & x\ge 0\\ -x & x<0 \end{cases} \]
\end{minipage}
\end{demo}
Bracket styles by name: \texttt{pmatrix}=( ), \texttt{bmatrix}=[ ],
\texttt{Bmatrix}=\{ \}, \texttt{vmatrix}=$|\;|$, \texttt{Vmatrix}=$\|\;\|$,
\texttt{matrix}=none.

\subsection{Spacing, text, braces and boxed answers}
Spaces: thin \verb|\,|, medium \verb|\:|, thick \verb|\;|, negative
\verb|\!|, and \verb|\quad|/\verb|\qquad| for big gaps. Upright words in
math: \verb|\text{if } x>0|. Annotate with braces:
\verb|\underbrace{x^2+2x+1}_{\text{square}}| $=\underbrace{x^2+2x+1}_{\text{square}}$
and \verb|\overbrace{a+b}^{n}| $=\overbrace{a+b}^{n}$. Box a final result:
\verb|\boxed{x=\frac{-b\pm\sqrt{\Delta}}{2a}}| $=\boxed{x=\tfrac{-b\pm\sqrt{\Delta}}{2a}}$.

% ======================================================================
\section{Symbol Reference (glyph + code)}
Every entry shows the printed symbol and the command that makes it. Most
need \texttt{amsmath}/\texttt{amssymb} (already loaded here).

\subsection{Lowercase Greek}
\renewcommand{\arraystretch}{1.3}
\begin{tabular}{@{}llllllll@{}}
\toprule
\sym{\alpha} & \sym{\beta} & \sym{\gamma} & \sym{\delta}\\
\sym{\epsilon} & \sym{\varepsilon} & \sym{\zeta} & \sym{\eta}\\
\sym{\theta} & \sym{\vartheta} & \sym{\iota} & \sym{\kappa}\\
\sym{\lambda} & \sym{\mu} & \sym{\nu} & \sym{\xi}\\
\sym{\pi} & \sym{\varpi} & \sym{\rho} & \sym{\varrho}\\
\sym{\sigma} & \sym{\varsigma} & \sym{\tau} & \sym{\upsilon}\\
\sym{\phi} & \sym{\varphi} & \sym{\chi} & \sym{\psi}\\
\sym{\omega} & \sym{\digamma} & \sym{\varkappa} &\\
\bottomrule
\end{tabular}

\subsection{Uppercase Greek}
(Letters that look like Latin --- A, B, E, Z, H, I, K, M, N, O, P, T, X ---
are just typed as capitals.)
\begin{tabular}{@{}llllllll@{}}
\toprule
\sym{\Gamma} & \sym{\Delta} & \sym{\Theta} & \sym{\Lambda}\\
\sym{\Xi} & \sym{\Pi} & \sym{\Sigma} & \sym{\Upsilon}\\
\sym{\Phi} & \sym{\Psi} & \sym{\Omega} & \sym{\aleph}\\
\bottomrule
\end{tabular}
Slanted variants (with \texttt{amsmath}): \verb|\varGamma \varDelta|
$\ldots$ $\varGamma\,\varDelta\,\varTheta\,\varPhi\,\varOmega$.

\subsection{Binary operators}
\begin{tabular}{@{}llllllll@{}}
\toprule
\sym{\pm} & \sym{\mp} & \sym{\times} & \sym{\div}\\
\sym{\cdot} & \sym{\ast} & \sym{\star} & \sym{\circ}\\
\sym{\bullet} & \sym{\oplus} & \sym{\ominus} & \sym{\otimes}\\
\sym{\oslash} & \sym{\odot} & \sym{\cup} & \sym{\cap}\\
\sym{\uplus} & \sym{\sqcup} & \sym{\sqcap} & \sym{\wedge}\\
\sym{\vee} & \sym{\setminus} & \sym{\wr} & \sym{\diamond}\\
\sym{\bigtriangleup} & \sym{\bigtriangledown} & \sym{\triangleleft} & \sym{\triangleright}\\
\sym{\dagger} & \sym{\ddagger} & \sym{\amalg} & \sym{\boxplus}\\
\sym{\boxtimes} & \sym{\boxdot} & \sym{\ltimes} & \sym{\rtimes}\\
\bottomrule
\end{tabular}

\subsection{Relations}
\begin{tabular}{@{}llllllll@{}}
\toprule
\sym{\leq} & \sym{\geq} & \sym{\neq} & \sym{\equiv}\\
\sym{\approx} & \sym{\cong} & \sym{\simeq} & \sym{\sim}\\
\sym{\propto} & \sym{\ll} & \sym{\gg} & \sym{\asymp}\\
\sym{\doteq} & \sym{\prec} & \sym{\succ} & \sym{\preceq}\\
\sym{\succeq} & \sym{\subset} & \sym{\supset} & \sym{\subseteq}\\
\sym{\supseteq} & \sym{\sqsubseteq} & \sym{\sqsupseteq} & \sym{\in}\\
\sym{\ni} & \sym{\notin} & \sym{\vdash} & \sym{\dashv}\\
\sym{\models} & \sym{\perp} & \sym{\mid} & \sym{\parallel}\\
\sym{\bowtie} & \sym{\smile} & \sym{\frown} & \sym{\therefore}\\
\sym{\because} & \sym{\triangleq} & \sym{\lesssim} & \sym{\gtrsim}\\
\bottomrule
\end{tabular}

\subsection{Negated relations}
\begin{tabular}{@{}llllllll@{}}
\toprule
\sym{\neq} & \sym{\notin} & \sym{\nleq} & \sym{\ngeq}\\
\sym{\nless} & \sym{\ngtr} & \sym{\nsim} & \sym{\ncong}\\
\sym{\nsubseteq} & \sym{\nsupseteq} & \sym{\subsetneq} & \sym{\supsetneq}\\
\sym{\nprec} & \sym{\nsucc} & \sym{\nmid} & \sym{\nparallel}\\
\sym{\nvdash} & \sym{\nvDash} & \sym{\nleftarrow} & \sym{\nrightarrow}\\
\bottomrule
\end{tabular}

\subsection{Arrows}
\begin{tabular}{@{}llll@{}}
\toprule
\sym{\leftarrow} & \sym{\rightarrow}\\
\sym{\leftrightarrow} & \sym{\Leftarrow}\\
\sym{\Rightarrow} & \sym{\Leftrightarrow}\\
\sym{\longleftarrow} & \sym{\longrightarrow}\\
\sym{\longleftrightarrow} & \sym{\Longrightarrow}\\
\sym{\Longleftrightarrow} & \sym{\mapsto}\\
\sym{\longmapsto} & \sym{\hookleftarrow}\\
\sym{\hookrightarrow} & \sym{\uparrow}\\
\sym{\downarrow} & \sym{\updownarrow}\\
\sym{\Uparrow} & \sym{\Downarrow}\\
\sym{\nearrow} & \sym{\searrow}\\
\sym{\swarrow} & \sym{\nwarrow}\\
\sym{\rightleftharpoons} & \sym{\leftharpoonup}\\
\sym{\rightharpoonup} & \sym{\twoheadrightarrow}\\
\sym{\rightarrowtail} & \sym{\leadsto}\\
\sym{\rightrightarrows} & \sym{\rightleftarrows}\\
\bottomrule
\end{tabular}
\quad Labelled arrows (need \texttt{amsmath}):
\verb|\xrightarrow{f}| $\xrightarrow{f}$,
\verb|\xleftarrow{g}| $\xleftarrow{g}$.

\subsection{Big (variable-size) operators}
\begin{tabular}{@{}llllllll@{}}
\toprule
\sym{\sum} & \sym{\prod} & \sym{\coprod} & \sym{\int}\\
\sym{\oint} & \sym{\iint} & \sym{\iiint} & \sym{\bigcup}\\
\sym{\bigcap} & \sym{\bigsqcup} & \sym{\bigvee} & \sym{\bigwedge}\\
\sym{\bigoplus} & \sym{\bigotimes} & \sym{\bigodot} & \sym{\biguplus}\\
\bottomrule
\end{tabular}

\subsection{Delimiters}
\begin{tabular}{@{}llllllll@{}}
\toprule
\sym{\langle} & \sym{\rangle} & \sym{\lfloor} & \sym{\rfloor}\\
\sym{\lceil} & \sym{\rceil} & \sym{\|} & \sym{\vert}\\
\sym{\Vert} & \sym{\backslash} & \sym{\ulcorner} & \sym{\urcorner}\\
\bottomrule
\end{tabular}

\subsection{Dots and accents}
\begin{tabular}{@{}llllllll@{}}
\toprule
\sym{\dots} & \sym{\cdots} & \sym{\vdots} & \sym{\ddots}\\
\sym{\hat{a}} & \sym{\bar{a}} & \sym{\tilde{a}} & \sym{\vec{a}}\\
\sym{\dot{a}} & \sym{\ddot{a}} & \sym{\acute{a}} & \sym{\grave{a}}\\
\sym{\breve{a}} & \sym{\check{a}} & \sym{\dddot{a}} & \sym{\widehat{abc}}\\
\sym{\widetilde{abc}} & \sym{\overline{abc}} & \sym{\underline{abc}} & \sym{\overrightarrow{abc}}\\
\bottomrule
\end{tabular}

\subsection{Miscellaneous symbols}
\begin{tabular}{@{}llllllll@{}}
\toprule
\sym{\infty} & \sym{\partial} & \sym{\nabla} & \sym{\emptyset}\\
\sym{\varnothing} & \sym{\forall} & \sym{\exists} & \sym{\nexists}\\
\sym{\neg} & \sym{\top} & \sym{\bot} & \sym{\angle}\\
\sym{\measuredangle} & \sym{\triangle} & \sym{\square} & \sym{\blacksquare}\\
\sym{\blacktriangle} & \sym{\lozenge} & \sym{\blacklozenge} & \sym{\eth}\\
\sym{\clubsuit} & \sym{\diamondsuit} & \sym{\heartsuit} & \sym{\spadesuit}\\
\sym{\flat} & \sym{\natural} & \sym{\sharp} & \sym{\Re}\\
\sym{\Im} & \sym{\wp} & \sym{\ell} & \sym{\hbar}\\
\sym{\imath} & \sym{\jmath} & \sym{\aleph} & \sym{\complement}\\
\sym{\prime} & \sym{\surd} & \sym{\checkmark} & \sym{\maltese}\\
\sym{\circledR} & \sym{\circledS} & \sym{\bigstar} & \sym{\mho}\\
\bottomrule
\end{tabular}

\subsection{Special alphabets}
\begin{tabular}{@{}ll@{}}
\toprule
Output & Code \\\midrule
$\mathbb{R}\ \mathbb{N}\ \mathbb{Z}\ \mathbb{Q}\ \mathbb{C}$
  & \verb|\mathbb{R} \mathbb{N} ...|\quad(blackboard, \texttt{amssymb})\\
$\mathcal{A}\ \mathcal{B}\ \mathcal{L}\ \mathcal{F}$
  & \verb|\mathcal{A} \mathcal{L} ...|\quad(calligraphic)\\
$\mathscr{A}\ \mathscr{B}\ \mathscr{L}\ \mathscr{F}$
  & \verb|\mathscr{A} ...|\quad(script, \texttt{mathrsfs})\\
$\mathfrak{A}\ \mathfrak{g}\ \mathfrak{p}\ \mathfrak{m}$
  & \verb|\mathfrak{A} ...|\quad(fraktur, \texttt{amssymb})\\
$\mathbf{A}\ \mathrm{A}\ \mathsf{A}\ \mathtt{A}\ \bm{A}$
  & \verb|\mathbf \mathrm \mathsf \mathtt \bm|\\
\bottomrule
\end{tabular}

\subsection{Named functions / operators}
Type these as commands so they print upright with correct spacing:
\verb|\sin \cos \tan \cot \sec \csc \arcsin \arccos \arctan \sinh|
\verb|\cosh \tanh \coth \exp \log \ln \lg \lim \limsup \liminf \sup|
\verb|\inf \max \min \arg \deg \det \dim \gcd \hom \ker \Pr|. For a
\emph{custom} operator with limits use
\verb|\DeclareMathOperator*{\argmax}{arg\,max}| in the preamble.

\subsection{Text-mode symbols (outside math)}
\begin{tabular}{@{}llllllll@{}}
\toprule
\tsym{\dag} & \tsym{\ddag} & \tsym{\S} & \tsym{\P}\\
\tsym{\copyright} & \tsym{\textregistered} & \tsym{\texttrademark} & \tsym{\pounds}\\
\tsym{\textbullet} & \tsym{\textdegree} & \tsym{\textbackslash} & \tsym{\ldots}\\
\tsym{\&} & \tsym{\%} & \tsym{\#} & \tsym{\_}\\
\bottomrule
\end{tabular}
\quad Check/cross via \texttt{pifont}: \verb|\ding{51}|\,\ding{51}\
\verb|\ding{55}|\,\ding{55}.

% ======================================================================
\section{Tables}
\subsection{The basics}
A \texttt{tabular} needs a \emph{column spec}: \texttt{l}/\texttt{c}/\texttt{r}
(left/centre/right), \verb=|= for a vertical rule, \verb|p{3cm}| for a fixed
wrapping width. Inside: \verb|&| separates columns, \verb|\\| ends a row,
\verb|\hline| draws a full horizontal rule, \verb|\cline{2-3}| a partial one.
\begin{demo}
\begin{minipage}[t]{0.5\textwidth}
\begin{lstlisting}[style=tex]
\begin{tabular}{|l|c|r|}
  \hline
  Left & Mid & Right\\ \hline
  a & b & c\\ \hline
\end{tabular}
\end{lstlisting}
\end{minipage}\hfill
\begin{minipage}[t]{0.45\textwidth}\centering
\begin{tabular}{|l|c|r|}
  \hline Left & Mid & Right\\ \hline a & b & c\\ \hline
\end{tabular}
\end{minipage}
\end{demo}

\subsection{multirow, multicolumn, colour and math in cells}
\verb|\multirow{N}{*}{txt}| spans \texttt{N} rows; \verb|\multicolumn{N}{|c|}{txt}|
spans \texttt{N} columns (and \emph{re-declares} the borders). Colour with
\verb|\rowcolor| (must be the \textbf{first} token of the row),
\verb|\cellcolor|, \verb|\columncolor|. Put maths in a cell with \verb|$...$|.
\begin{demo}
\begin{minipage}[t]{0.54\textwidth}
\begin{lstlisting}[style=tex]
\renewcommand{\arraystretch}{1.4}
\begin{tabular}{|c|c|c|}
\hline
\rowcolor{gray!30}
\multicolumn{2}{|c|}{\textbf{Header}} & C\\
\hline
\multirow{2}{*}{Merged}
  & \cellcolor{yellow!30}$x^2$ & $s>0$\\
\cline{2-3}
  & $\dfrac{1}{s}$ & $s>0$\\
\hline
\end{tabular}
\end{lstlisting}
\end{minipage}\hfill
\begin{minipage}[t]{0.42\textwidth}\centering
\renewcommand{\arraystretch}{1.4}
\begin{tabular}{|c|c|c|}
\hline
\rowcolor{gray!30}
\multicolumn{2}{|c|}{\textbf{Header}} & C\\ \hline
\multirow{2}{*}{Merged} & \cellcolor{yellow!30}$x^2$ & $s>0$\\ \cline{2-3}
 & $\dfrac{1}{s}$ & $s>0$\\ \hline
\end{tabular}
\end{minipage}
\end{demo}
Taller rows: \verb|\renewcommand{\arraystretch}{1.5}|. Thicker rules:
\verb|\setlength{\arrayrulewidth}{1pt}| or use \texttt{booktabs}
(\verb|\toprule \midrule \bottomrule|, no vertical rules --- the modern look).
Line breaks inside a cell: \verb|\makecell{line 1\\ line 2}| (\texttt{makecell}).

\subsection{Auto-width tables (tabularx)}
When a column must stretch to fill a set total width, use \texttt{tabularx}
and the \texttt{X} column type:
\begin{demo}
\ttfamily\small
\textbackslash begin\{tabularx\}\{0.8\textbackslash textwidth\}\{|l|X|\}
\end{demo}
\begin{tabularx}{0.8\textwidth}{|l|X|}
\hline Key & This X-column wraps automatically to fill the declared width so
the table always spans exactly 80\% of the text width.\\ \hline
\end{tabularx}

% ======================================================================
\section{Figures and Graphics}
\subsection{One figure}
\begin{lstlisting}[style=tex]
\begin{figure}[htbp]          % h=here t=top b=bottom p=floats-page
  \centering
  \includegraphics[width=0.5\textwidth]{example-image-a}
  \caption{A single figure.}  % \caption BEFORE \label
  \label{fig:single}
\end{figure}
\end{lstlisting}
\begin{figure}[H]
  \centering
  \includegraphics[width=0.45\textwidth]{example-image-a}
  \caption{A single figure (rendered here with \texttt{[H]} from \texttt{float}).}
  \label{fig:single}
\end{figure}
Scaling keys for \verb|\includegraphics|: \verb|width=|, \verb|height=|,
\verb|scale=0.5|, \verb|angle=90| (rotate), \verb|keepaspectratio|. Force an
exact position with \verb|[H]| (needs \texttt{float}); use \verb|[htbp!]| to
strongly suggest.

\subsection{Side-by-side subfigures}
\texttt{subcaption} gives each panel its own (a),(b) sub-label.
\begin{lstlisting}[style=tex]
\begin{figure}[htbp]\centering
  \begin{subfigure}[b]{0.3\textwidth}
    \includegraphics[width=\linewidth]{example-image-a}
    \caption{First}\label{fig:a}
  \end{subfigure}\hfill
  \begin{subfigure}[b]{0.3\textwidth}
    \includegraphics[width=\linewidth]{example-image-b}
    \caption{Second}\label{fig:b}
  \end{subfigure}
  \caption{Two panels.}\label{fig:pair}
\end{figure}
\end{lstlisting}
\begin{figure}[H]\centering
  \begin{subfigure}[b]{0.3\textwidth}\centering
    \includegraphics[width=\linewidth]{example-image-a}\caption{First}\label{fig:a}
  \end{subfigure}\hfill
  \begin{subfigure}[b]{0.3\textwidth}\centering
    \includegraphics[width=\linewidth]{example-image-b}\caption{Second}\label{fig:b}
  \end{subfigure}
  \caption{Two panels (rendered).}\label{fig:pair}
\end{figure}
\verb|\ref{fig:a}| gives the full number ``\ref{fig:a}''; \verb|\subref{fig:a}|
gives just the letter. For panels \emph{without} sub-captions, use
\verb|\begin{minipage}{0.48\textwidth}...\end{minipage}| instead.

\subsection{Wrapped and rotated figures}
\begin{lstlisting}[style=tex]
\begin{wrapfigure}{r}{0.35\textwidth}  % r=right side, l=left side
  \centering
  \includegraphics[width=\linewidth]{example-image}
  \caption{Text flows around me.}\label{fig:wrap}
\end{wrapfigure}
% ---- rotation ----
\includegraphics[angle=90]{example-image}  % rotate the image
\rotatebox{45}{slanted text}               % rotate anything (graphicx)
\begin{sidewaysfigure}...\end{sidewaysfigure}  % 90-deg whole float (rotating)
\end{lstlisting}
Rotated text sample: \rotatebox{20}{\fbox{tilted 20\textdegree}}.

% ======================================================================
\section{Cross-References and Hyperlinks}
Label anything right after its \verb|\caption| (or in an equation), then
point at it:
\begin{tabular}{@{}ll@{\hspace{2em}}l@{}}
\toprule
Command & Prints & Note\\\midrule
\verb|\ref{k}|     & the bare number   & figures, tables, sections\\
\verb|\eqref{k}|   & (number)          & equations --- adds parentheses\\
\verb|\pageref{k}| & the page number   & ``see p.~\pageref{fig:single}''\\
\verb|\autoref{k}| & ``Figure~3''      & adds the word automatically\\
\verb|\cite{k}|    & [1]               & links to a \verb|\bibitem|\\
\bottomrule
\end{tabular}\\[4pt]
External links (\texttt{hyperref}): clickable text
\verb|\href{https://example.com}{visit site}| $\to$
\href{https://example.com}{visit site}; bare URL
\verb|\url{https://example.com}|. Always put a \verb|~| (non-breaking space)
before a reference so the number never wraps alone: \verb|Figure~\ref{fig:pair}|.
Configure link colours once with \verb|\hypersetup{colorlinks=true,...}|.

% ======================================================================
\section{Code Listings}
The \texttt{listings} package typesets source code verbatim with optional
syntax colouring and line numbers.
\begin{lstlisting}[style=py, caption={Inline Python}, label={lst:py}]
def rc_response(R, C, V, times):
    """One-line docstring."""
    tau = R * C
    return [V * (1 - 2.718**(-t / tau)) for t in times]
\end{lstlisting}
Inline snippet: \verb|\lstinline!print(x)!|. From an external file:
\verb|\lstinputlisting[style=py]{script.py}|. Define a reusable style once
with \verb|\lstdefinestyle{py}{language=Python, numbers=left, frame=single,|
\verb|keywordstyle=\color{blue}, ...}|. For publication-grade output the
\texttt{minted} package is prettier but needs \verb|-shell-escape| + Python.

% ======================================================================
\section{Theorems, Proofs and Coloured Boxes}
Declare environments in the preamble with \texttt{amsthm}:
\verb|\newtheorem{theorem}{Theorem}[section]| numbers them ``1.1, 1.2''.
\begin{definition}
A \emph{prime} is an integer $>1$ divisible only by $1$ and itself.
\end{definition}
\begin{theorem}[Pythagoras]\label{thm:pyth}
In a right triangle, $a^2 + b^2 = c^2$.
\end{theorem}
\begin{proof}
Standard area argument. \qedhere
\end{proof}
The \texttt{proof} environment adds the italic ``Proof.'' lead-in and the
$\square$. A \emph{coloured, titled} theorem box via \texttt{tcolorbox}:
\begin{keyresult}{Weak-field orbit}{weak}
$\left(\dfrac{v}{c}\right)^2 = \dfrac{GM}{rc^2} = \dfrac{r_s}{2r}.$
\end{keyresult}
Reference them like anything else: Theorem~\ref{thm:pyth}.

% ======================================================================
\section{Algorithms / Pseudocode}
Using \texttt{algorithm2e} (loaded with \texttt{[ruled,vlined,linesnumbered]}):
\begin{algorithm}[H]
\SetAlgoLined
\KwIn{array $A$ of length $n$}
\KwOut{maximum element}
$m \gets A[1]$\;
\For{$i \gets 2$ \KwTo $n$}{
  \If{$A[i] > m$}{$m \gets A[i]$\;}
}
\Return $m$\;
\caption{Find the maximum}
\end{algorithm}
Key commands: \verb|\KwIn \KwOut \For \While \If \ElseIf \Else \Return|,
end each statement line with \verb|\;|. The alternative package pair is
\texttt{algorithm} + \texttt{algpseudocode} (\verb|\State \For \EndFor|).

% ======================================================================
\section{Footnotes and Bibliography}
A footnote\footnote{This text drops to the bottom of the page and
auto-numbers.} is placed with \verb|\footnote{...}| immediately after the
word. The fastest bibliography for a timed exam is the manual
\texttt{thebibliography} environment (no BibTeX run needed):
\begin{lstlisting}[style=tex]
\begin{thebibliography}{9}   % {9} = widest label is one digit
  \bibitem{kreyszig} E. Kreyszig, \textit{Advanced Engineering
    Mathematics}, 10th ed. Wiley, 2011.
  \bibitem{ogata} K. Ogata, \textit{Modern Control Engineering},
    5th ed. Prentice Hall, 2010.
\end{thebibliography}
\end{lstlisting}
Cite in the text with \verb|\cite{kreyszig}|, which prints [1] and links to
the entry.

% ======================================================================
\section{Predicted Tricky Variations}
\begin{enumerate}[label=\textbf{\arabic*.},leftmargin=*]
\item \textbf{Multi-column aligned maths} --- when two equation pairs sit
  side by side, \texttt{align} spaces them oddly; use \texttt{alignat}, where
  you control the gaps:
  \verb|\begin{alignat}{2} a&=b &\quad c&=d \end{alignat}|.
\item \textbf{Numbered sub-equations} (1a),(1b): wrap an \texttt{align} in
  \verb|\begin{subequations}...\end{subequations}|.
\item \textbf{Numbered derivation, one number} --- keep \verb|&=| aligned and
  end each intermediate line with \verb|\nonumber| so only the last is numbered.
\item \textbf{Clean scientific notation / units} --- \texttt{siunitx}:
  \verb|\num{1.496e11}| $\to 1.496\times10^{11}$ and \verb|\SI{9.8}{m/s^2}|.
\item \textbf{Coloured full-row header} --- \verb|\rowcolor| must be the FIRST
  token of the row, even before \verb|\multicolumn|.
\item \textbf{Landscape float} --- \texttt{sidewaystable}/\texttt{sidewaysfigure}
  (\texttt{rotating}) or the \verb|{landscape}| environment (\texttt{lscape}).
\item \textbf{Custom numbered coloured theorem} --- \texttt{tcolorbox} with
  \verb|\newtcbtheorem| (as the ``Key Result'' box above).
\item \textbf{Tensor / pre-indices} --- \verb|{}_a^b X| for pre-scripts, and a
  phantom space \verb|\ | (or \verb|\vphantom|) to stop indices stacking.
\item \textbf{Continued fractions} --- \verb|\cfrac{1}{1+\cfrac{1}{1+x}}|.
\item \textbf{Highlight/box a final answer} --- \verb|\boxed{...}| or
  \verb|\colorbox{yellow}{$...$}|.
\end{enumerate}

% ======================================================================
\section{Common Errors and Instant Fixes}
\renewcommand{\arraystretch}{1.3}
\begin{tabular}{@{}p{0.42\textwidth}p{0.54\textwidth}@{}}
\toprule
Message / symptom & Fix\\\midrule
\texttt{Missing \$ inserted} & A maths symbol used in text --- wrap it in \verb|$...$|.\\
\texttt{Undefined control sequence} & Typo, or the package that defines it isn't loaded.\\
\texttt{! Misplaced alignment tab \&} & A bare \verb|&| outside a table/align --- escape it \verb|\&|.\\
\verb|\ref| shows ``??'' & Compile again (labels resolve on the 2nd pass).\\
Float ``lost'' / on wrong page & Use \verb|[htbp]|, add \verb|!|, or \verb|[H]| (\texttt{float}).\\
Table too wide & Shrink columns, use \texttt{tabularx}, or \verb|\small| / \verb|\resizebox|.\\
Fractions collide in a table & \verb|\renewcommand{\arraystretch}{1.5}|.\\
\texttt{Option clash for package xcolor} & Load \texttt{xcolor} (with options) \emph{before} \texttt{tcolorbox}/\texttt{tikz}.\\
Accented letter breaks & Load \verb|\usepackage[T1]{fontenc}| and \verb|[utf8]{inputenc}|.\\
\bottomrule
\end{tabular}

% ======================================================================
\section{Exam Speed Strategy}
\begin{enumerate}[leftmargin=*]
\item Paste a preamble first; add packages only as each feature appears.
\item Build in reading order: preamble $\to$ title $\to$ sections/lists $\to$
  equations $\to$ tables $\to$ figures $\to$ code $\to$ references.
\item \textbf{Compile every few environments} so any error is easy to locate.
\item \textbf{Copy--paste--tweak} repetitive blocks (table rows, subfigures,
  \verb|\bibitem|s) --- never retype them.
\item Keep this compiled PDF open on a second window for instant symbol lookup.
\item Leave line breaks, exact float positions and page breaks for last ---
  they are usually \emph{not} graded, but successful compilation always is.
\end{enumerate}

\vfill
\begin{center}\small\itshape
End of reference --- compile twice, print, and keep beside you.
\end{center}

\end{document}
